% Generated draft: requires mathematical and source review.
\documentclass[11pt]{article}
\usepackage[margin=25mm]{geometry}
\usepackage{fontspec}
\setmainfont{Latin Modern Roman}
\usepackage{amsmath,amssymb,mathtools,unicode-math}
\setmathfont{Latin Modern Math}
\usepackage{enumitem,needspace,xurl,hyperref}
\setlength{\parindent}{0pt}
\setlength{\parskip}{5pt}
\setlength{\emergencystretch}{3em}
\newcommand{\N}{\mathbb{N}}
\newcommand{\Z}{\mathbb{Z}}
\newcommand{\Q}{\mathbb{Q}}
\newcommand{\R}{\mathbb{R}}
\newcommand{\C}{\mathbb{C}}
\newcommand{\F}{\mathbb{F}}
\newcommand{\A}{\mathbb{A}}
\newcommand{\PP}{\mathbb{P}}
\newcommand{\E}{\mathbb{E}}
\newcommand{\eps}{\varepsilon}
\newcommand{\dd}{\,\mathrm{d}}
\begin{document}
\section*{UnsolvedMath ID: 30004223}
\textbf{Source problem code:} catalogue entry
\subsection{Definitions and mathematical statement}
Let \(d\in\Z\) with \(d\geq4\). Under the usual convex-polytope convention, a \(d\)-polytope is a convex polytope of affine dimension \(d\). It is simplicial if every proper nonempty face is a simplex. For such a polytope \(P\), define
\[
f(P)=(f_0(P),\ldots,f_{d-1}(P)),\qquad
f_i(P)=\#\{F:F\text{ is an }i\text{-dimensional face of }P\}.
\]
Here the empty face and \(P\) itself are omitted. A finite regular CW \((d-1)\)-sphere is a finite regular CW complex \(K\) whose underlying space is homeomorphic to \(S^{d-1}\); regularity means that each characteristic map restricts to a homeomorphism from the closed ball onto the closure of its cell.

For either \(X=\partial P\), with its proper-face cell structure, or such a complex \(X=K\), and for \(S=\{s_1<\cdots<s_k\}\subseteq\{0,\ldots,d-1\}\), define
\[
f_S(X)=\#\bigl\{(e_1,\ldots,e_k):\dim e_j=s_j,\ e_j\subseteq\overline{e_{j+1}}\text{ for }1\leq j<k\bigr\},
\qquad f_{\varnothing}(X)=1.
\]
The flag \(f\)-vector is \(\operatorname{flag}f(X)=(f_S(X))_{S\subseteq\{0,\ldots,d-1\}}\). Write \(\operatorname{flag}f(P)=\operatorname{flag}f(\partial P)\).

Define the collections
\[
\begin{aligned}
\mathcal F_d^{\mathrm{simp}}&=\{f(P):P\text{ is a simplicial }d\text{-polytope}\},\\
\mathcal F_d^{\mathrm{poly}}&=\{f(P):P\text{ is a }d\text{-polytope}\},\\
\mathcal H_d^{\mathrm{poly}}&=\{\operatorname{flag}f(P):P\text{ is a }d\text{-polytope}\},\\
\mathcal H_{d-1}^{\mathrm{CW}}&=\{\operatorname{flag}f(K):K\text{ is a finite regular CW }(d-1)\text{-sphere}\}.
\end{aligned}
\]
\textbf{Source research question.} For each integer \(d\geq4\), identify qualitative differences between \(\mathcal F_d^{\mathrm{simp}}\) and \(\mathcal F_d^{\mathrm{poly}}\), and between \(\mathcal H_d^{\mathrm{poly}}\) and \(\mathcal H_{d-1}^{\mathrm{CW}}\).

\textbf{Unresolved formulation qualifications.} The source does not define ``qualitative differences'' or prescribe a property to prove, a type of separation, or a classification criterion. Thus this is an exploratory comparison problem, not a uniquely specified conjecture; in particular, it does not explicitly assert strict inequality of either pair of collections in every dimension. The definitions above make the objects and coordinates explicit using conventional convex polytopes, finite CW complexes, and ordinary topological spheres. The supplied statement does not expressly specify convexity, finiteness, or a stronger category of sphere, so these conventions require human confirmation against the source.
\subsection{Short English statement}
For every integer \(d\geq4\), compare the possible face-count vectors of simplicial and general \(d\)-polytopes, and compare the possible flag face-count vectors of polytope boundaries and regular CW \((d-1)\)-spheres. Flag counts record chains of faces or cells of prescribed dimensions. The requested ``qualitative differences'' are unspecified, so the source supplies an open-ended research question rather than a definite theorem to establish.
\subsection{Sources}
\begin{itemize}
\item Local UnsolvedMath registry, permanent ID \(30004223\), source problem code \texttt{catalogue entry}, title \emph{Qualitative Differences Among Polytope and Sphere Face Vectors}, in the Oberwolfach Reports Open Problems collection. The supplied record gives proposed year \(2019\) and names no individual proposer.
\item Supplied source citation: \emph{Geometric, Algebraic, and Topological Combinatorics} (2020), Oberwolfach Reports, pp.\ 2395--2472, DOI \texttt{10.4171/owr/2019/39}. \url{https://doi.org/10.4171/owr/2019/39}. The supplied background describes Oberwolfach Reports as documenting mathematical workshops and their open-problem discussions.
\item Supplied literature source, labeled ``Oberwolfach report'': \url{https://doi.org/10.4171/owr/2019/39}.
\item Supplied literature source, labeled ``Sjoeberg and Ziegler'': \url{https://arxiv.org/abs/1711.01864}. No title or specific supporting passage is supplied for this reference.
\item Supplied literature assessment, checked on \texttt{2026-08-22}: the record labels the problem open and states that the \(g\)-theorem completely describes simplicial \(f\)-vectors, that qualitative separation from general and flag \(f\)-vector sets remains unresolved, and that the \(4\)-polytope \(f\)-vector set is still not characterized. This is supplied dated curation derived from the OWR material, not an independently verified literature assessment. The supplied background expressly requests independent verification of time-sensitive status and cited claims.
\item The registry labels the statement \texttt{corrected\_verified} and explains that notation, typography, and sentence structure were repaired while preserving mathematical content. This is a supplied editorial qualification, not independent mathematical verification of this draft.
\end{itemize}
\end{document}
